% LREC 2026 Example; 
% LREC Is now using templates similar to the ACL ones. 
\documentclass[10pt, a4paper]{article}

\usepackage[review]{lrec2026} % this is the new style
% the 'review' option anonymizes the paper following submission guideline
% the 'final' option produces the camera ready version (non anonymized)
% default version is 'final', so use review option for submission
\usepackage{tabularx}
\usepackage{amsfonts}

\title{Hybrid Privacy-Preserving Text Anonymization via Risk-Aware Differential Privacy and k-Anonymity}

\name{Yamaç Eren Ay, Ibrahim Baroud} 

\address{Technische Universität Berlin \\
         yamac.ay@campus.tu-berlin.de, ibrahim.baroud@tu-berlin.de\\
}

% Each paper must include an abstract of 150 to 200 words in 9 pt with interlinear spacing of 10 pt. The heading Abstract should be centered, 10 pt bold.

\abstract{In this article, we investigate a hybrid approach called k-DP-MLM for text anonymization that combines Differential Privacy (DP) with empirical k-anonymity measures. The goal is to extend context-aware token-level rewriting (achieved by DP-MLM \cite{meisenbacher_dp-mlm_2024}) by k-anonymity guided selective rewriting (achieved by PETRE \cite{manzanares-salor_enhancing_2025}), instead of applying global noise to the entire sentence. Upon testing the new approach against both baselines, we observe that k-DP-MLM is able to anonymize a record up to an arbitrary privacy level (determined by the noise hyperparameter), while preserving text utility and quality. \Keywords{K-Anonymization, Differential Privacy, Token-level Rewriting} }

\begin{document}

\maketitleabstract

% \section{My List of Points}

% I suggest this way:

% We start with a mock anonymization importance paragraph, saying that it's important to keep the data private, and make it still utility-preserving.

% There are lots of papers discussing it. We see that there are multiple different definitions of privacy. For example, the simplest methods we see online are the custom NER-based maskers, which remove the identifiers from the records directly. Usually, they're used for removing the direct PIIs, but indirect PIIs can be removed as well by annotating and training a custom NER model on that (Baroud). Those methods usually have a classification threshold, which defines the minimum confidence level to be able to classify some entity with a label.

% The other types of anonymity definitions are k-anonymity. In textual data, they're difficult to achieve, partly because we have to consider textual ambiguity and all. PETRE is a good option for those who want to mask iteratively until the probability k-anonymity measures are satisfied. It uses XAI and explainability scores to measure the re-id risk of individual tokens.

% The final types of anonymity definitions we talk about is the differential privacy. It's literally a new paradigm shift compared to k-anonymity, saying that we need to make changes to the data so that the value changes are indistinguishable at any adversarial input etc. I don't know the exact definition but sth like that. There are two types of anonymizers which work like that. First ones are the full text rewriters, which take the text, encode it in the parametric memory, then perturb the data by the given epsilon (privacy budget), then decode back to obtain the privatized text. DP-Prompt, DPBart work like that. In addition, there are obfuscation methods for example DP-Paraphrase, which are prompted to anonymize the given text, so they rewrite the text as well. But the most important type of methods we are gonna talk about are the token-level rewriters, for example DP-MLM. DP-MLM doesn't really invent a real big thing, they already have something like that in the literature, but it uses static embeddings to replace the word and ignores the whole context. DP-MLM uses BERT embeddings to get the full context and act upon that, anonymizing individual tokens one by one.

% We find out that the DP-MLM is not optimal yet. Applying global, uniform noise to all tokens is hard to fine-tune in practice: The noise has to be set high enough to replace the most disclosive tokens revealing the identity of the individual, but setting it too high may potentially lead to utility losses.

% This way, we take inspiration from PETRE and integrate three ideas into our DP-MLM variant. Firstly, make noise proportional to the re-identification risk of individual tokens, so that we can optimize the privacy-utility trade-off while preserving the formal guarantees by distributing the epsilon smartly towards the risky tokens. Secondly, start from the most risky tokens, and process them in descending order, so that the probability of re-identification sinks the fastest over the course of token-wise anonymization. Thirdly, define a stopping condition which is evaluated at each anonymization step and tells whether the rank of the current document is not below the desired rank or not. We process each record until we reach the desired rank k.

% To be more specific, our proposed methods are k-DP-MLM with risk-aware noising + sorting: Use SHAP or greedy-based explanation scores to estimate the marginal contribution of individual tokens to the overall re-identification (where SHAP is combinatorial and greedy is meant to be faster in terms of worst case complexity). The weights we obtain for each token is inverted (reflecting "the higher the risk, the lower the epsilon, thus more likely to anonymize" logic) and scaled (with the given epsilon as the mean of the new distribution). For testing the effect of sorting, we define the ablation variants: k-DP-MLM with risk-aware noising but no sorting, vanilla k-DP-MLM, vanilla DP-MLM. We expect to have a more utility-preserving DP-MLM in k-DP-MLM. And the risk scores should guide us to find the most risky tokens to save time to anonymize, thus improving both utility and time complexity as well. And the sorting should possibly help us get faster to the k-anonymity.

% Funnily enough, our well-proposed full variant is identical to PETRE but with rewriting capability rather than masking.

\section{Introduction}

Text anonymization is a fundamental prerequisite for responsibly sharing and processing textual data in domains such as law, healthcare, and social media, where sensitive personal information is often embedded in free-form text. The main challenge lies in balancing privacy protection against data utility: overly aggressive anonymization can render text unusable for downstream tasks, while insufficient anonymization exposes individuals to re-identification risks. As a result, a wide range of anonymization techniques has been proposed, each grounded in different definitions of privacy and threat models.

A common class of practical anonymization methods relies on Named Entity Recognition (NER) and related span detection techniques. These approaches identify explicit personal identifiers, such as names or locations, and replace them with placeholders or generalized labels. While originally designed for removing direct identifiers, supervised NER models can also be extended and trained to detect indirect or contextual identifiers if enough annotation samples are provided in the training dataset. These PII detection methods usually define the classification threshold as a hyperparameter, which determines how aggressively tokens are treated as sensitive and can be fine-tuned for individual cases. Although such methods are efficient and easy to deploy, they lack formal privacy guarantees and remain vulnerable to inference attacks that exploit residual information.

Another line of work is grounded in k-anonymity, which aims to ensure that each record in a released dataset is indistinguishable from at least $k-1$ other records. Translating this notion to text is challenging due to linguistic variability and contextual ambiguity. PETRE \cite{manzanares-salor_enhancing_2025} addresses this challenge by operationalizing probabilistic k-anonymity for textual data, using explainability methods to estimate token-level re-identification risk and iteratively masking the most disclosive tokens until the desired anonymity level is reached. PETRE is originally run on de-identified records to solely handle the most disclosive indirect identifiers. While effective in reducing empirical re-identification risk, PETRE relies on deterministic masking and does not provide formal, tunable privacy guarantees.

A fundamentally different paradigm is offered by Differential Privacy (DP) \cite{dwork_differential_2006}, which limits how much the output distribution can change when a single record is modified. Existing DP-based anonymizers can be broadly categorized into full-text rewriting methods, such as DP-BART \cite{igamberdiev_dp-bart_2023} and DP-Prompt \cite{utpala_locally_2023}, which encode the entire input text into latent representations, perturb them according to a privacy budget $\epsilon$, and decode a privatized version of the text. While these methods provide formal guarantees, they often incur high computational costs and may introduce semantic drift.

More fine-grained approaches operate at the token level. DP-MLM \cite{meisenbacher_dp-mlm_2024} applies differential privacy to masked language model–based rewriting, anonymizing text by replacing individual tokens in context. By leveraging bidirectional contextual embeddings, DP-MLM improves coherence compared to earlier embedding-based substitution methods that ignore context. However, DP-MLM applies a global, uniform noise budget across all tokens, which makes it difficult to fine-tune in practice. The privacy budget must be large enough to perturb the most disclosive tokens, but doing so uniformly can lead to unnecessary modifications of low-risk tokens, degrading utility and increasing runtime.

In this work, we address these limitations by integrating ideas from empirical risk-based anonymization into the DP-MLM framework. Specifically, we propose a hybrid approach that combines formal differential privacy guarantees with token-level re-identification risk estimation inspired by PETRE. Our goal is to achieve stronger privacy–utility trade-offs, improved runtime efficiency, and a principled stopping criterion aligned with probabilistic k-anonymity.

\section{Preliminaries}\label{sec:preliminaries}

This section formalizes the threat model, introduces text re-identification (TRI) attacks, and defines the empirical notion of k-anonymity used throughout this work. It also summarizes the differential privacy definition that underpins our rewriting mechanism. These preliminaries establish the adversarial setting under which anonymization mechanisms are evaluated and motivate the design choices of k-DP-MLM.

\subsection{Text Re-Identification Attacker}

We consider a \emph{text re-identification} (TRI) attacker whose goal is to link an anonymized text record to its originating individual. Formally, let $\mathcal{D} = \{(x_i, u_i)\}_{i=1}^N$ denote a dataset of textual records $x_i$ associated with user identities $u_i$. Given an anonymized version $\tilde{x}$ of a target record $x$, the attacker outputs a ranked list of candidate identities from a finite set $\mathcal{U}$. The attacker is modeled as a classifier $f_\theta$ trained to estimate the posterior probability
\[
p(u \mid \tilde{x}),
\]
where $\theta$ denotes learned parameters. In practice, $f_\theta$ is implemented as a transformer-based text classifier trained on auxiliary background knowledge. We assume the attacker is \emph{static}: the model is trained once and remains fixed during anonymization and evaluation.

\subsection{Empirical k-Anonymity via Rank}

Given the attacker’s posterior distribution, we define the \emph{rank} of the true identity $u^\ast$ as its position in the attacker’s sorted prediction list:
\[
\mathrm{rank}(\tilde{x}) = \mathrm{rank}_{u^\ast} \big( p(u \mid \tilde{x}) \big).
\]

A record is said to satisfy \emph{empirical k-anonymity} if the true identity cannot be distinguished from at least $k-1$ other candidates under the attacker model, i.e.,
\[
\mathrm{rank}(\tilde{x}) \geq k.
\]

This rank-based definition provides an operational notion of probabilistic k-anonymity for free-form text and is commonly used in text anonymization and re-identification studies. Unlike classical k-anonymity, which relies on equivalence classes over structured attributes, this formulation naturally accounts for linguistic variability and contextual information.

\subsection{Background Knowledge and Static Adversary Assumption}

We assume the attacker has access to extensive auxiliary background knowledge about users, such as prior writings or contextual metadata, which may be de-identified but is not fully anonymized. This background knowledge is used to train the TRI model before any anonymization is applied to the target records.

Crucially, the attacker is assumed to be \emph{static}: it does not adapt its parameters in response to anonymized outputs. This assumption is standard in scalable text anonymization evaluations, as retraining the attacker after each anonymization step would be computationally prohibitive. Moreover, static attackers capture realistic scenarios in which adversaries accumulate background knowledge prior to targeting specific anonymized texts.

\subsection{Token-Level Re-Identification Risk}

To guide selective anonymization, we associate each token in a text with a \emph{re-identification risk score} that estimates its contribution to successful re-identification. These scores are computed with respect to the TRI attacker using explainability methods, such as SHAP-based attribution \cite{lundberg_unified_2017}, which quantify the marginal impact of individual tokens on the attacker’s prediction.

We denote the risk score of token $t_j$ as $r_j \geq 0$, where larger values indicate higher disclosiveness. Risk scores may reflect direct identifiers (e.g., names) as well as indirect or contextual identifiers. In practice, direct identifiers detected by prior de-identification can be assigned arbitrarily large risk or an effectively zero privacy budget, ensuring they are always anonymized.

These token-level risk estimates form the basis for risk-aware noise allocation and processing order in k-DP-MLM.

\subsection{Transformers}

Transformers \cite{vaswani_attention_2017} replace recurrent sequence modeling with self-attention, allowing each token to attend to all others in the input. This yields contextual token representations that capture long-range dependencies and can be computed efficiently in parallel. We use transformer encoders, which produce a representation for every token in the sequence.

\subsection{Masked Language Models and BERT}

Masked language models (MLMs) train transformer encoders to predict a token that has been replaced by a mask symbol, producing a distribution over the vocabulary for that position. BERT \cite{devlin_bert_2019} is a standard encoder-only transformer pretrained with the MLM objective and then fine-tuned for downstream classification or regression tasks. This setup reuses pretrained linguistic representations and requires only task-specific supervision. We use MLMs as conditional predictors when selecting token replacements in DP-MLM and related methods.

\subsection{Differential Privacy Preliminaries}\label{sec:dp-prelim}

Differential Privacy (DP) \cite{dwork_differential_2006} is defined for a randomized mechanism $\mathcal{M}$ as follows: for any neighboring datasets $D$ and $D'$ that differ in one record and any measurable set of outputs $S$,
\[
\Pr[\mathcal{M}(D) \in S] \le e^{\epsilon} \Pr[\mathcal{M}(D') \in S].
\]
The probabilities are taken over the internal randomness of $\mathcal{M}$. The bound limits how much the output distribution can shift when a single record changes. The parameter $\epsilon$ is the privacy budget that controls the strength of the guarantee. We adopt the same formulation and use it to parameterize the rewriting mechanism, following DP-MLM \cite{meisenbacher_dp-mlm_2024}.

\subsection{Overview on DP-MLM}

DP-MLM \cite{meisenbacher_dp-mlm_2024} is a token-level DP-based text anonymization method that rewrites a sentence one token at a time. For each step, it masks a token, runs a pretrained masked language model (MLM), and obtains a distribution over candidate replacements conditioned on the full context. The replacement is sampled with temperature scaling derived from the privacy budget, so lower $\epsilon$ implies higher temperature and a higher chance of changing the token. This sampling implements the exponential mechanism \cite{mcsherry_mechanism_2007} under bounded logits, yielding a differentially private selection rule. The process iterates over all tokens to produce the anonymized text.

\section{k-DP-MLM}

We introduce k-DP-MLM, a family of extensions to DP-MLM that incorporate three key ideas.

\subsection{Risk-aware noise allocation}

Instead of applying a single privacy budget to all tokens, we allocate noise proportionally to token-level re-identification risk. In DP-MLM, a global $\epsilon$ is used across the sequence, so low-risk tokens are perturbed as aggressively as high-risk ones. We estimate token risks with an explainability method over a static TRI attacker (SHAP by default) and map the raw scores to a normalized scale because the values are numerically close to zero. Let
\[
s_i = \frac{r_{\max} - r_i}{r_{\max} - r_{\min}}
\]
be the normalized risk score for token $i$, where larger $s_i$ indicates higher disclosure risk. We assign a token-specific privacy budget $\epsilon_i$ that is co-ordered with $s_i$ and scaled to preserve the global mean, $\mathbb{E}_i[\epsilon_i] = \epsilon$. A simple choice is the softmax normalization:
\begin{equation}
\epsilon_i = \epsilon \cdot n \cdot \frac{\exp (s_i / \tau) }{\sum_{j=1}^{n} \exp (s_j / \tau)},
\end{equation}
where $\tau > 0$ controls the sharpness of the allocation, with smaller values concentrating budget adjustments on the highest-risk tokens.

When all tokens are rewritten under this allocation, the overall DP budget remains $\epsilon$ in expectation. As in PETRE-style pipelines, explicit de-identification is still recommended; direct identifiers can be treated as near-zero $\epsilon$ (or equivalently very large risk), while the proposed allocation focuses on residual indirect identifiers.

\subsection{Risk-based processing order}

Secondly, we process tokens in descending order of estimated risk. By prioritizing the most disclosive tokens, the probability of successful re-identification decreases rapidly during early anonymization steps. This ordering not only improves privacy more efficiently but also reduces unnecessary rewriting of low-risk tokens, leading to better utility and lower computational cost.

\subsection{k-anonymity–driven stopping condition}

Thirdly, we introduce a dynamic stopping criterion based on empirical k-anonymity. After each anonymization step, we evaluate the rank of the current document under a TRI attacker. Anonymization continues until the rank of the true identity is no better than k, at which point the record is considered sufficiently anonymized. This allows each record to be processed only as much as necessary, further improving efficiency and utility.

\subsection{Variants and ablations}

Concretely, we evaluate the following variants:
\begin{itemize}
    \item \textbf{k-DP-MLM (risk-aware + sorting).} Risk-proportional noise allocation with descending-risk token order and k-based stopping. Prior de-identification is recommended.
    \item \textbf{k-DP-MLM (risk-aware, no sorting).} Risk-proportional noise allocation with arbitrary token order and k-based stopping. Prior de-identification is recommended.
    \item \textbf{DP-MLM (risk-aware).} Risk-proportional noise allocation (with arbitrary token order). Prior de-identification is not necessary.
    \item \textbf{DP-MLM (uniform).} Uniform noise applied to all tokens without risk awareness or stopping.
\end{itemize}

We hypothesize that risk-aware noise allocation improves utility by focusing perturbations on truly sensitive tokens, while sorting accelerates convergence to k-anonymity and reduces runtime. Notably, the full k-DP-MLM variant is structurally analogous to PETRE, differing primarily in that it performs context-aware rewriting under differential privacy, rather than deterministic masking.

For PETRE comparisons, we use a one-shot risk ordering computed once before anonymization so that runtime assumptions match the DP-MLM pipeline.

\section{Data}

\subsection{Datasets}

We use two datasets: TAB and PersonalReddit. TAB \cite{pilan_text_2022} provides a controlled setting for text re-identification with approximately 1.5k record samples that include text, metadata, and entity annotations, which enables PII detection training and evaluation. For privacy and utility assessment, we use the test split with 127 records from 127 distinct individuals. We define a lightweight downstream task for TAB, such as predicting macro-region or year of the court case, to quantify utility without adding new annotation work.

PersonalReddit \cite{staab_beyond_2024} consists of persona-grounded dialogues with demographic attributes that support both privacy and utility evaluation. The train and test files contain 525 answers from 40 distinct individuals in total, which provides enough identity repetition for re-identification experiments. We use the demographic attributes as downstream labels for utility evaluation.

\subsection{Data Preparation}

We construct proxy background knowledge for training TRI models to avoid overfitting to the target records. This background knowledge is ideally de-identified and theoretically unbounded, but can also be derived from the original records when external data is unavailable. Models are trained on this auxiliary information and validated and tested on the de-identified target records. This setup reflects an adversary who acquires contextual knowledge prior to targeting specific instances and introduces lexical variation that reduces memorization.

\section{Experimental Setup}

\subsection{TRI Model Training and Selection}

To select suitable text re-identification (TRI) models for each dataset, we performed a hyperparameter grid search over learning rates $ \{ 2 \times 10^{-5}, 5 \times 10^{-5}, 1 \times 10^{-4}, 2 \times 10^{-4} \} $. The resulting configurations were chosen based on validation performance and training stability.

We use two TRI pipelines per dataset. The baseline pipeline trains on original background knowledge and evaluates on original records to produce baseline accuracy and a reference rank distribution. The de-identified pipeline trains on Presidio-de-identified background knowledge and evaluates on de-identified records for rank evaluation and token-risk estimation. Presidio \footnote{\url{https://microsoft.github.io/presidio/}} is an open-source NER-based PII detection and masking tool.

For TAB, the optimal setting used 20 training epochs with pretrained initialization, a batch size of 16, and a learning rate of $5 \times 10^{-5}$. 

For PersonalReddit, the best configuration used 15 epochs, pretraining on the corpus, a batch size of 32, and a learning rate of $5 \times 10^{-5}$.

\begin{table}[!htp]
\centering
\small
\setlength{\tabcolsep}{4pt}
\begin{tabularx}{\columnwidth}{|l|X|X|l|}
\hline
\textbf{Dataset} & \textbf{Background Knowledge (Train)} & \textbf{Records (Eval)} & \textbf{Accuracy} \\
\hline
TAB & Original & Original & 0.8492 \\
TAB & De-identified using Presidio & De-identified using Presidio & 0.7778 \\
Reddit & Original & Original & 1.0000 \\
Reddit & De-identified using Presidio & De-identified using Presidio & 1.0000 \\
\hline
\end{tabularx}
\caption{TRI pipeline accuracy for original and de-identified settings on TAB and PersonalReddit.}
\end{table}

\subsection{Privacy Evaluation}

Privacy experiments follow the PETRE-style protocol: a text re-identification model is trained to link records to identities using the selected textual feature. The trained re-identifier is then used to measure empirical re-identification risk on anonymized texts via rank. This setup aligns with the static attacker assumption and keeps evaluation scalable.

\subsection{Utility Evaluation}

Utility is measured with downstream classification on the original datasets. We train a task model on the original texts and evaluate it on both original and anonymized outputs to compute relative utility loss. For multi-class settings, we use TF-IDF vectors and logistic regression as a simple, reproducible baseline \cite{alemerien_sentiment_2024}.

\subsection{Text Similarity}

As a complementary signal, we compute the distance between original and anonymized texts to characterize the privacy--utility trade-off. We report BLEU and cosine similarity as lightweight baselines, and include BERTScore and MoverScore to capture contextual similarity.

% \section{Results}

% We report privacy, utility, and runtime results for all variants and datasets.

% \begin{table}[!htp]
% \centering
% \small
% \setlength{\tabcolsep}{4pt}
% \begin{tabularx}{\columnwidth}{|l|X|X|X|}
% \hline
% \textbf{Method} & \textbf{Privacy} & \textbf{Utility} & \textbf{Runtime} \\
% \hline
%  &  &  &  \\
% \hline
% \end{tabularx}
% \caption{Placeholder results table to be filled with privacy, utility, and runtime metrics.}
% \label{tab:results-overview}
% \end{table}


% To select the best text re-identification models for each dataset, we conducted hyperparameter grid search with learning rates $2 \times 10^{-5}, 5 \times 10^{-5}, 10^{-4}, 2 \times 10^{-4}$. The optimal hyperparameters for the training:

% \begin{itemize}
%     \item TAB dataset: Epochs = 20, Pretraining, Batch size = 16, Learning rate = $5 \times 10^{-5}$. Reached $78.57\%$ accuracy.
%     \item Reddit dataset: Epochs = 15, Pretraining, Batch size = 32, Learning rate = $5 \times 10^{-5}$. Reached $95\%$ accuracy.
% \end{itemize}

% To enable fair comparison across different anonymization strategies and to facilitate a clearer analysis of computational complexity and runtime performance, we introduce a simplified variant of PETRE. In the original PETRE framework, re-identification risk is recomputed for each interesting \footnote{Interesting means we exclude stop words and special tokens masking direct identifiers} token at every anonymization step. In contrast, our simplified formulation computes token-level risk scores only once, prior to anonymization, and subsequently processes tokens in descending order of these precomputed scores. This modification aligns PETRE more closely with DP-MLM, which similarly performs a single global risk evaluation at the outset, thereby allowing for a more meaningful comparison between the two methods.\\

% Importantly, these experiments are not intended to demonstrate limitations or weaknesses of PETRE. Rather, the goal is to normalize methodological differences so that both approaches operate under comparable assumptions and computational constraints. We assume that explainability-based risk estimators, such as SHAP, already account for higher-order and combinatorial interactions among tokens when estimating marginal contributions. Under this assumption, the relative ordering of tokens by risk is expected to remain largely stable across anonymization iterations.\\

% Beyond algorithmic considerations, we also revisit the training and evaluation protocol of the token re-identification (TRI) models. In standard machine learning practice, training, validation, and test sets are kept strictly disjoint to prevent information leakage and overly optimistic performance estimates. However, we observed that the training strategy adopted in the original PETRE implementation, where TRI models are trained and validated exclusively on the same records with no difference whether de-identified or original, leads to severe overfitting. In particular, the absence of an effective stopping criterion allows models to achieve near-perfect training accuracy without necessarily learning generalizable re-identification patterns.\\

% To address this issue, we constructed proxy background knowledge, which is ideally de-identified and in fact theoretically unbounded but also can also be generated from the original records, and used this information for training while validating and testing on the de-identified target records. This approach reflects a more realistic scenario in which an adversary acquires contextual knowledge before targeting specific records, and it introduces lexical and structural variation that mitigates memorization effects.\\

% Finally, we frame our privacy evaluation within a static attacker model. We assume an adversary with access to extensive auxiliary information about users, such as demographic attributes or socioeconomic indicators, who seeks to re-identify anonymized texts by exploiting leaked indirect identifiers. This setting reflects a realistic and widely studied threat model, while also offering favorable scalability properties: the attacker’s models are trained once on the original data and subsequently evaluated on privatized outputs. In contrast, adaptive attack scenarios would require retraining after each anonymization step, significantly increasing computational cost. For these reasons, we restrict our analysis to the static attacker setting in this work.

% \section{Introduction}

% The widespread adoption of data-driven natural language processing (NLP) systems has led to an unprecedented reliance on large-scale textual datasets. Many of these datasets contain sensitive or personally identifiable information (PII), raising concerns about privacy, data protection, and responsible use. Legal frameworks such as the General Data Protection Regulation (GDPR) impose increasingly strict constraints on how textual data may be processed, shared, and reused, particularly in settings where the risk of re-identification cannot be ruled out. In this context, text anonymization has become a fundamental requirement for modern language technology, situated at the intersection of legal compliance, ethical responsibility, and technical feasibility.\\

% Early text anonymization methods typically rely on surface-level transformations, most commonly the detection and masking of named entities. Although such approaches can effectively remove explicit identifiers, they often fail to prevent contextual disclosure, where sensitive information remains inferable from surrounding text. To address these limitations, recent work has explored neural rewriting methods that modify text at the token or sentence level. These approaches aim to reduce disclosure risks while preserving semantic coherence, but they introduce a well-known tension between providing meaningful privacy protection and maintaining text utility.\\

% Differential Privacy (DP) offers a principled way to address this tension by providing formal, quantifiable privacy guarantees that are independent of an adversary’s background knowledge. Applied to language data, DP-based mechanisms introduce controlled noise into model outputs to limit the influence of individual records. DP-MLM \cite{meisenbacher_dp-mlm_2024} follows this paradigm by applying differential privacy to masked language model-based rewriting, enabling context-aware anonymization at the token level. While DP-MLM achieves strong theoretical guarantees, its reliance on global noise injection can lead to unnecessary modifications of non-sensitive tokens, resulting in avoidable degradation of text quality.\\

% Along a different line of work, k-anonymity-inspired approaches focus on reducing re-identification risk by ensuring that sensitive elements are indistinguishable among at least $k$ similar instances. PETRE \cite{manzanares-salor_enhancing_2025} exemplifies this strategy by using empirical k-anonymity measures to guide selective token rewriting. By targeting only those tokens that contribute most to re-identification risk, PETRE is able to preserve more of the original text structure and meaning. However, this selectivity comes at the cost of lacking formal and tunable privacy guarantees comparable to those provided by differential privacy.\\

% In this paper, we introduce \textit{k-DP-MLM}, a hybrid anonymization approach that combines differential privacy with k-anonymity-guided selective rewriting. The proposed method integrates the formal privacy guarantees of DP-MLM with the targeted rewriting strategy of PETRE, aiming to limit unnecessary perturbations while retaining explicit control over privacy levels through a single noise hyperparameter. By unifying these two perspectives, k-DP-MLM seeks to balance theoretical privacy protection with practical utility preservation.\\

% We evaluate k-DP-MLM against both DP-MLM and PETRE across multiple privacy and utility dimensions. The results show that k-DP-MLM can anonymize textual records to arbitrary privacy levels while consistently maintaining higher text quality and downstream utility than either baseline. These findings highlight the potential of hybrid privacy mechanisms as a practical and scalable solution for legally compliant text anonymization in real-world NLP pipelines.


% \section{How to Produce the \texttt{.pdf}}
% \label{sec:append-how-prod}

% In order to generate a PDF file out of the LaTeX file herein, when citing language resources, the following steps need to be performed:

% \begin{enumerate}
% \item \texttt{xelatex your\_paper\_name.tex}
% \item \texttt{bibtex your\_paper\_name.aux}
% \item \texttt{bibtex languageresource.aux}    *NEW*
% \item \texttt{xelatex your\_paper\_name.tex}
% \item \texttt{xelatex your\_paper\_name.tex}
% \end{enumerate}

% From 2026 we are using a sans-serif font TeX Gyre Heros, so you must
% install it if you do not have it.  For monospaced font, we use TeX
% Gyre Cursor.\footnote{They are available from CTAN in the
%   \href{https://ctan.org/pkg/tex-gyre-heros}{tex-gyre-heros} and
%   \href{https://ctan.org/pkg/tex-gyre-cursor}{tex-gyre-cursor} packages
%     or in TeXLive as \texttt{tex-gyre} and \texttt{tex-cursor}.}  To
%   compile, you can use either \textbf{\texttt{pdfLaTeX}},
%   \textbf{\texttt{XeLaTeX}} or \textbf{\texttt{LuaLaTeX}}.

% \section{Final Paper}

% Each final paper should be submitted online. The fully justified text should be formatted according to LREC 2026 style as indicated for the Full Paper submission.

% Final papers can be five (5) to nine (9)  pages in length, including figures (plus more pages if needed for references, ethical consideration, limitations, conflict-of-interest, as well as data and code availability statements).  
% Length does not depend on the presentation mode (oral or poster).

% \begin{itemize}
% \item The paper is in A4-size format, that is 21 x 29.7 cm.
% \item The text height is 24.7 cm and the text width 16.0 cm in two
%   columns separated by a 0.6 cm space.
% \item The font for the entire paper should be Tex Gyre Heros, a
%   sans-serif font based Helvetica, with TeX Gyre Cursor for the
%   monospaced font.
% \item The main text should be 10 pt with an interlinear spacing of 11
%   pt.
% \item The use of \texttt{lrec2026.sty} will ensure the good formatting.
% \end{itemize}



% \subsection{General Instructions for the Final Paper}

% The unprotected PDF files will appear in the online proceedings directly as received. \textbf{Do not print the page number}.

% \section{Page Numbering}

% \textbf{Please do not include page numbers in your Paper.} The definitive page numbering of papers published in the proceedings will be decided by the Editorial Committee.

% \section{Headings / Level 1 Headings} 

% Level 1 Headings should be capitalised in the same way as the main title, and centered within the column. The font used is 12 pt bold. There should also be a space of 12 pt between the title and the preceding section and 3 pt between the title and the following text.

% \subsection{Level 2 Headings}

% The format of Level 2 Headings is the same as for Level 1 Headings, with the font  11 pt, and the heading is justified to the left of the column. There should also be a space of 6 pt between the title and the preceding section and 3 pt between the title and the following text.

% \subsubsection{Level 3 Headings}
% \label{level3H}

% The format of Level 3 Headings is the same as Level 2, except that the font is  10 pt, and there should be no space left between the heading and the text as in~\ref{level3H}. There should also be a space of 6 pt between the title and the preceding section and 3 pt between the title and the following text.

% \section{Citing References in the Text}

% \subsection{Bibliographical References}


% \begin{table}
% \centering
% \begin{tabular}{lll}
% \hline
% \textbf{Output} & \textbf{natbib command} & \textbf{Old command}\\
% \hline
% \citep{Eco:1990} & \verb|\citep| & \verb|\cite| \\
% \citealp{Eco:1990} & \verb|\citealp| & no equivalent \\
% \citet{Eco:1990} & \verb|\citet| & \verb|\newcite| \\
% \citeyearpar{Eco:1990} & \verb|\citeyearpar| & \verb|\shortcite| \\
% \hline
% \end{tabular}
% \caption{\label{citation-guide} Citation commands supported by the style file. The style is based on the natbib package and supports all natbib citation commands. It also supports commands defined in previous style files for compatibility.}
% \end{table}

% Table~\ref{citation-guide} shows the syntax supported by the style files. We encourage you to use the natbib styles.
% You can use the command \verb|\citet| (cite in text) to get ``author (year)'' citations, like this citation to a paper by \citet{CastorPollux-92}. You can use the command \verb|\citep| (cite in parentheses) to get ``(author, year)'' citations \citep{CastorPollux-92}. You can use the command \verb|\citealp| (alternative cite without parentheses) to get ``author, year'' citations, which is useful for using citations within parentheses (e.g. \citealp{CastorPollux-92}).

% When several authors are cited, those references should be separated with a semicolon: \cite{Martin-90,CastorPollux-92}. When the reference has more than three authors, only cite the name of the first author followed by ``et.\ al.'', e.g. \citet{Superman-Batman-Catwoman-Spiderman-00}.

% \subsection{Language Resource References}

% \subsubsection{When Citing Language Resources}

% As you may know, LREC introduced a separate section on Language
% Resources citation to enhance the value of such assets. When citing
% language resources, we recommend to proceed in the same way as for
% bibliographical references. Please make sure to compile your Language
% Resources stored as a \texttt{.bib} file \textbf{separately}. This
% produces the required \texttt{.aux} and \texttt{.bbl} files.  See
% Section~\ref{sec:append-how-prod} for details on how to produce this
% with bibtex.

% A language resource should be cited inline as, for example,
% \citetlanguageresource{Speecon}, or \citetlanguageresource{EMILLE} or
% parenthetically \citeplanguageresource{ItalWordNet}.  All online
% language resources should have a persistent identifier (pid), following the \href{https://doi.org/10.15497/RDA00040}{Tromsø recommendations for citation of research data in linguistics}.

% Please add either an ISLRN (International Standard Language Resource
% Number) in the \texttt{islrn} field or some other url in the
% \texttt{pid} field.  This will be hyperlinked and shown in the
% bibliography.  We give an example in the entry for \texttt{ItalWordNet} which you can see in \path{languageresource.bib}.
  



  
% \section{Figures \& Tables}

% \subsection{Figures}

% All figures should be centred and clearly distinguishable. They should never be drawn by hand, and the lines must be very dark in order to ensure a high-quality printed version. Figures should be numbered in the text, and have a caption in 10 pt underneath. A space must be left between each figure and its respective caption. 

% Example of a figure:

% \begin{figure}[!ht]
% \begin{center}
% \includegraphics[width=\columnwidth]{home-palacio.jpg}
% \caption{The caption of the figure.}
% \label{fig.1}
% \end{center}
% \end{figure}

% Figure and caption should always appear together on the same page. Large figures can be centered, using a full page.

% \subsection{Tables}

% The instructions for tables are the same as for figures.

% \begin{table}[!ht]
% \begin{center}
% \begin{tabularx}{\columnwidth}{|l|X|}

%       \hline
%       Level&Tools\\
%       \hline
%       Morphology & Pitrat Analyser\\
%       \hline
%       Syntax & LFG Analyser (C-Structure)\\
%       \hline
%      Semantics & LFG F-Structures + Sowa's\\
%      & Conceptual Graphs\\
%       \hline

% \end{tabularx}
% \caption{The caption of the table}
%  \end{center}
% \end{table}

% \section{Footnotes}

% Footnotes are indicated within the text by a number in superscript\footnote{Footnotes should be in  9 pt, and appear at the bottom of the same page as their corresponding number. Footnotes should also be separated from the rest of the text by a 5 cm long horizontal line.}.

% \section{Copyrights}

% The Language Resources and Evaluation Conference (LREC) Proceedings are published by the European Language Resources Association (ELRA). They are available online from the conference website.

% ELRA's policy is to acquire copyright for all LREC contributions. In assigning your copyright, you are not forfeiting your right to use your contribution elsewhere. This you may do without seeking permission and is subject only to normal acknowledgment to the LREC proceedings. The LREC Proceedings are licensed under CC-BY-NC, the Creative Commons Attribution-Non-Commercial 4.0 International License.


% \section{Conclusion}

% Your submission of a finalized contribution for inclusion in the LREC Proceedings automatically assigns the above copyright to ELRA.

% \section{Acknowledgements}

% Place all acknowledgments (including those concerning research grants and funding) in a separate section at the end of the paper.

% \section{Optional Supplementary Materials: Appendices, Software, and Data}

% \subsection{Appendices}

% \textbf{Appendices or supplementary material will be allowed ONLY in the final, camera-ready version, but not during submission, as papers should be reviewed without the need to refer to any supplementary materials.}

% Each \textbf{camera ready} submission can be accompanied by an appendix, up to ten (10) pages long, usually being included in a main PDF paper file, one \texttt{.tgz} or \texttt{.zip} archive containing software, and one \texttt{.tgz} or \texttt{.zip} archive containing data.

% We encourage the submission of these supplementary materials to improve the reproducibility of results and to enable authors to provide additional information that does not fit in the paper. For example, preprocessing decisions, model parameters, feature templates, lengthy proofs or derivations, pseudocode, sample system inputs/outputs, and other details necessary for the exact replication of the work described in the paper can be put into the appendix. However, the paper submissions must remain fully self-contained, as these supplementary materials are optional. If the pseudo-code or derivations, or model specifications are an essential part of the contribution, or if they are important for the technical correctness of the work, they should be a part of the main paper and not appear in the appendix.

% Appendices are material that can be read and include lemmas, formulas, proofs, and tables that are not critical to the reading and understanding of the paper, as in \href{https://acl-org.github.io/ACLPUB/formatting.html#appendices}{*ACLPUB}. \textbf{It is  highly recommended that the appendices should come after the references}; the main text and appendices should be contained in a `single' manuscript file, without being separately maintained. Letter them in sequence and provide an informative title: \textit{Appendix A. Title of Appendix}


% \subsection{Extra space for ethical considerations and limitations}

% Please note that extra space is allowed after the last page for an ethics/broader impact statement and a discussion of limitations. At submission time, if you need extra space for these sections, it should be placed after the conclusion so that it is possible to rapidly check that the rest of the paper still fits in 8 pages maximum. Ethical considerations sections, limitations, acknowledgments, and references do not count against these limits.

% \section{Providing References}

% \subsection{Bibliographical References} 

% Bibliographical references should be listed in alphabetical order at the end of the paper. The title of the section, ``Bibliographical References'', should be a Level 1 Heading. The first line of each bibliographical reference should be justified to the left of the column, and the rest of the entry should be indented by 0.35 cm.

% The examples provided in Section~\ref{sec:reference} (some of which are fictitious references) illustrate the basic format required for papers in conference proceedings, books, journal articles, PhD theses, and books chapters.

% \subsection{Language Resource References}

% Language resource references should be listed in alphabetical order at the end of the paper.

% \nocite{Eco:1990,Martin-90,Chercheur-94,CastorPollux-92,bs-2570-manual,bs-2570-techreport,chomsky-73,hoel-71-portion,hoel-71-whole,singer-whole,Jespersen:1922,Superman-Batman-Catwoman-Spiderman-00}
\section{Bibliographical References}\label{sec:reference}

\bibliographystyle{lrec2026-natbib}
\bibliography{lrec2026-example}

\section{Language Resource References}
\label{lr:ref}
\bibliographystylelanguageresource{lrec2026-natbib}

\end{document}

%%% Local Variables:
%%% mode: latex
%%% TeX-master: t
%%% End:
